\section{6G Service Orchestration}
\subsection{Preserving the Budget after Wireless Delivery}
For 6G mobile edge services, a short wireless transmission time is only one part of timely execution. A delivered request may still be blocked until its new service requirements become usable. Equation~\eqref{eq:remaining-budget} identifies this activation wait within the same send-to-completion budget. The measured-wait replays show how its timing, together with scheduling duration, translates into modeled completion differences under shared radio inputs and a common scheduler.

This supports treating intent-activation delay as part of the service QoS budget. Its value is clearest during transitions, when one interpretation can release many waiting requests. The image-service results show why both transition-window and full-trajectory completion matter: Jev completes more requests during updates than either generative comparator, while its full-run completion remains close to DeepSeek. These measurements connect interpretation speed to useful results returned after wireless delivery.

\subsection{Separating Intent Changes from Network Dynamics}
Mobility and changing edge load require repeated resource decisions even when the service intent remains unchanged. The integration results support caching the interpreted contract and using current numerical observations for those decisions. Serving-cell changes can alter the transport cost of a placement without changing its permission; a newly restricted location instead requires a new contract. Distinguishing these events avoids repeated language interpretation on the numerical scheduling path and retains adaptation to network and computing state. Fast decision models are most useful at first interpretation and subsequent contract updates.

\subsection{Choosing the Interpreter and Its Interface}
Interpreter selection should consider the output interface together with the model. Qwen's direct attributes improve correctness and completion despite a longer response, making that interface the relevant local comparison for Jev. The result favors explicit contract fields when a compact catalog identifier imposes a difficult mapping task. Evaluating the stronger interface retains Jev's latency advantage while comparing decisions of matched observed quality.

Latency and price favor different hosted models in these measurements. Across the cloud block's 21 calls per model, Jev reports US\$0.002125, DeepSeek US\$0.001262, and Gemini US\$0.009698. The image-service fees in Table~\ref{tab:application} also favor DeepSeek. Jev supplies the faster measured response; DeepSeek has the lower billed interpretation fee. Self-hosted Qwen has no per-call API invoice here, but uses a separate compute allocation. A deployment can weigh these fees against the completion benefit of faster activation.

\subsection{Study Scope and Next Steps}
The latency comparisons include the model, serving system, and transport path. Our replays exchange recorded model and scheduling waits together, measuring their effect on execution without attributing the delay to individual provider components. NR simulation supplies the wireless-access evidence for this study of 6G service orchestration; the image service adds actual transfers, queues, and OCR, with two logical workers on one host.

The experiments use a bounded contract catalog and authored transitions. Timing repeats reuse the same inputs, and the application adds two mobility/load traces; requests sharing a trajectory are dependent observations. Completion differences are descriptive, including the four-request Jev--DeepSeek gap in the image service. Varying mobility, link load, and distributed edge capacity independently would test how much interpretation waiting each communication and computing regime can tolerate. Independently sampled contracts would extend the semantic coverage.
